%!TeX root=../princesstop.tex
\chapter{The Magic}

\begin{figure}[t!]
\centering
\includegraphics[width=\linewidth]{filigree}
\end{figure}

\lettrine[lines=5]{W}{hen} Sara had passed the house next door she had seen Ram Dass closing the shutters, and caught her glimpse of this room also.

\zz
<It is a long time since I saw a nice place from the inside,> was the thought which crossed her mind.

\zz
There was the usual bright fire glowing in the grate, and the Indian gentleman was sitting before it. His head was resting in his hand, and he looked as lonely and unhappy as ever.

<Poor man!> said Sara. <I wonder what you are supposing.>

And this was what he was <supposing> at that very moment.

<Suppose,> he was thinking, <suppose—even if Carmichael traces the people to Moscow—the little girl they took from Madame Pascal's school in Paris is \textsc{not} the one we are in search of. Suppose she proves to be quite a different child. What steps shall I take next?>

When Sara went into the house she met Miss Minchin, who had come downstairs to scold the cook.

<Where have you wasted your time?> she demanded. <You have been out for hours.>

<It was so wet and muddy,> Sara answered, <it was hard to walk, because my shoes were so bad and slipped about.>

<Make no excuses,> said Miss Minchin, <and tell no falsehoods.>

Sara went in to the cook. The cook had received a severe lecture and was in a fearful temper as a result. She was only too rejoiced to have someone to vent her rage on, and Sara was a convenience, as usual.

<Why didn't you stay all night?> she snapped.

Sara laid her purchases on the table.

<Here are the things,> she said.

The cook looked them over, grumbling. She was in a very savage humour indeed.

<May I have something to eat?> Sara asked rather faintly.

<Tea's over and done with,> was the answer. <Did you expect me to keep it hot for you?>

Sara stood silent for a second.

<I had no dinner,> she said next, and her voice was quite low. She made it low because she was afraid it would tremble.

<There's some bread in the pantry,> said the cook. <That's all you'll get at this time of day.>

Sara went and found the bread. It was old and hard and dry. The cook was in too vicious a humour to give her anything to eat with it. It was always safe and easy to vent her spite on Sara. Really, it was hard for the child to climb the three long flights of stairs leading to her attic. She often found them long and steep when she was tired; but tonight it seemed as if she would never reach the top. Several times she was obliged to stop to rest. When she reached the top landing she was glad to see the glimmer of a light coming from under her door. That meant that Ermengarde had managed to creep up to pay her a visit. There was some comfort in that. It was better than to go into the room alone and find it empty and desolate. The mere presence of plump, comfortable Ermengarde, wrapped in her red shawl, would warm it a little.

Yes; there Ermengarde was when she opened the door. She was sitting in the middle of the bed, with her feet tucked safely under her. She had never become intimate with Melchisedec and his family, though they rather fascinated her. When she found herself alone in the attic she always preferred to sit on the bed until Sara arrived. She had, in fact, on this occasion had time to become rather nervous, because Melchisedec had appeared and sniffed about a good deal, and once had made her utter a repressed squeal by sitting up on his hind legs and, while he looked at her, sniffing pointedly in her direction.

<Oh, Sara,> she cried out, <I am glad you have come. Melchy \textsc{would} sniff about so. I tried to coax him to go back, but he wouldn't for such a long time. I like him, you know; but it does frighten me when he sniffs right at me. Do you think he ever \textsc{would} jump?>

<No,> answered Sara.

Ermengarde crawled forward on the bed to look at her.

<You \textsc{do} look tired, Sara,> she said; <you are quite pale.>

<I \textsc{am} tired,> said Sara, dropping on to the lopsided footstool. <Oh, there's Melchisedec, poor thing. He's come to ask for his supper.>

Melchisedec had come out of his hole as if he had been listening for her footstep. Sara was quite sure he knew it. He came forward with an affectionate, expectant expression as Sara put her hand in her pocket and turned it inside out, shaking her head.

<I'm very sorry,> she said. <I haven't one crumb left. Go home, Melchisedec, and tell your wife there was nothing in my pocket. I'm afraid I forgot because the cook and Miss Minchin were so cross.>

Melchisedec seemed to understand. He shuffled resignedly, if not contentedly, back to his home.

<I did not expect to see you tonight, Ermie,> Sara said. Ermengarde hugged herself in the red shawl.

<Miss Amelia has gone out to spend the night with her old aunt,> she explained. <No one else ever comes and looks into the bedrooms after we are in bed. I could stay here until morning if I wanted to.>

She pointed toward the table under the skylight. Sara had not looked toward it as she came in. A number of books were piled upon it. Ermengarde's gesture was a dejected one.

<Papa has sent me some more books, Sara,> she said. <There they are.>

Sara looked round and got up at once. She ran to the table, and picking up the top volume, turned over its leaves quickly. For the moment she forgot her discomforts.

<Ah,> she cried out, <how beautiful! Carlyle's French Revolution. I have \textsc{so} wanted to read that!>

<I haven't,> said Ermengarde. <And papa will be so cross if I don't. He'll expect me to know all about it when I go home for the holidays. What \textsc{shall} I do?>

Sara stopped turning over the leaves and looked at her with an excited flush on her cheeks.

<Look here,> she cried, <if you'll lend me these books, \textit{I'll} read them—and tell you everything that's in them afterward—and I'll tell it so that you will remember it, too.>

<Oh, goodness!> exclaimed Ermengarde. <Do you think you can?>

<I know I can,> Sara answered. <The little ones always remember what I tell them.>

<Sara,> said Ermengarde, hope gleaming in her round face, <if you'll do that, and make me remember, I'll—I'll give you anything.>

<I don't want you to give me anything,> said Sara. <I want your books—I want them!> And her eyes grew big, and her chest heaved.

<Take them, then,> said Ermengarde. <I wish I wanted them—but I don't. I'm not clever, and my father is, and he thinks I ought to be.>

Sara was opening one book after the other. <What are you going to tell your father?> she asked, a slight doubt dawning in her mind.

<Oh, he needn't know,> answered Ermengarde. <He'll think I've read them.>

Sara put down her book and shook her head slowly. <That's almost like telling lies,> she said. <And lies—well, you see, they are not only wicked—they're \textsc{vulgar}. Sometimes>—reflectively—<I've thought perhaps I might do something wicked—I might suddenly fly into a rage and kill Miss Minchin, you know, when she was ill-treating me—but I \textsc{couldn't} be vulgar. Why can't you tell your father \textit{I} read them?>

<He wants me to read them,> said Ermengarde, a little discouraged by this unexpected turn of affairs.

<He wants you to know what is in them,> said Sara. <And if I can tell it to you in an easy way and make you remember it, I should think he would like that.>

<He'll like it if I learn anything in \textsc{any} way,> said rueful Ermengarde. <You would if you were my father.>

<It's not your fault that\longdash> began Sara. She pulled herself up and stopped rather suddenly. She had been going to say, <It's not your fault that you are stupid.>

<That what?> Ermengarde asked.

<That you can't learn things quickly,> amended Sara. <If you can't, you can't. If I can—why, I can; that's all.>

She always felt very tender of Ermengarde, and tried not to let her feel too strongly the difference between being able to learn anything at once, and not being able to learn anything at all. As she looked at her plump face, one of her wise, old-fashioned thoughts came to her.

<Perhaps,> she said, <to be able to learn things quickly isn't everything. To be kind is worth a great deal to other people. If Miss Minchin knew everything on earth and was like what she is now, she'd still be a detestable thing, and everybody would hate her. Lots of clever people have done harm and have been wicked. Look at Robespierre\longdash>

She stopped and examined Ermengarde's countenance, which was beginning to look bewildered. <Don't you remember?> she demanded. <I told you about him not long ago. I believe you've forgotten.>

<Well, I don't remember \textsc{all} of it,> admitted Ermengarde.

<Well, you wait a minute,> said Sara, <and I'll take off my wet things and wrap myself in the coverlet and tell you over again.>

She took off her hat and coat and hung them on a nail against the wall, and she changed her wet shoes for an old pair of slippers. Then she jumped on the bed, and drawing the coverlet about her shoulders, sat with her arms round her knees. <Now, listen,> she said.

She plunged into the gory records of the French Revolution, and told such stories of it that Ermengarde's eyes grew round with alarm and she held her breath. But though she was rather terrified, there was a delightful thrill in listening, and she was not likely to forget Robespierre again, or to have any doubts about the Princesse de Lamballe.

<You know they put her head on a pike and danced round it,> Sara explained. <And she had beautiful floating blonde hair; and when I think of her, I never see her head on her body, but always on a pike, with those furious people dancing and howling.>

It was agreed that Mr St John was to be told the plan they had made, and for the present the books were to be left in the attic.

<Now let's tell each other things,> said Sara. <How are you getting on with your French lessons?>

<Ever so much better since the last time I came up here and you explained the conjugations. Miss Minchin could not understand why I did my exercises so well that first morning.>

Sara laughed a little and hugged her knees.

<She doesn't understand why Lottie is doing her sums so well,> she said; <but it is because she creeps up here, too, and I help her.> She glanced round the room. <The attic would be rather nice—if it wasn't so dreadful,> she said, laughing again. <It's a good place to pretend in.>

The truth was that Ermengarde did not know anything of the sometimes almost unbearable side of life in the attic and she had not a sufficiently vivid imagination to depict it for herself. On the rare occasions that she could reach Sara's room she only saw the side of it which was made exciting by things which were <pretended> and stories which were told. Her visits partook of the character of adventures; and though sometimes Sara looked rather pale, and it was not to be denied that she had grown very thin, her proud little spirit would not admit of complaints. She had never confessed that at times she was almost ravenous with hunger, as she was tonight. She was growing rapidly, and her constant walking and running about would have given her a keen appetite even if she had had abundant and regular meals of a much more nourishing nature than the unappetizing, inferior food snatched at such odd times as suited the kitchen convenience. She was growing used to a certain gnawing feeling in her young stomach.

<I suppose soldiers feel like this when they are on a long and weary march,> she often said to herself. She liked the sound of the phrase, <long and weary march.> It made her feel rather like a soldier. She had also a quaint sense of being a hostess in the attic.

<If I lived in a castle,> she argued, <and Ermengarde was the lady of another castle, and came to see me, with knights and squires and vassals riding with her, and pennons flying, when I heard the clarions sounding outside the drawbridge I should go down to receive her, and I should spread feasts in the banquet hall and call in minstrels to sing and play and relate romances. When she comes into the attic I can't spread feasts, but I can tell stories, and not let her know disagreeable things. I dare say poor chatelaines had to do that in time of famine, when their lands had been pillaged.> She was a proud, brave little chatelaine, and dispensed generously the one hospitality she could offer—the dreams she dreamed—the visions she saw—the imaginings which were her joy and comfort.

So, as they sat together, Ermengarde did not know that she was faint as well as ravenous, and that while she talked she now and then wondered if her hunger would let her sleep when she was left alone. She felt as if she had never been quite so hungry before.

<I wish I was as thin as you, Sara,> Ermengarde said suddenly. <I believe you are thinner than you used to be. Your eyes look so big, and look at the sharp little bones sticking out of your elbow!>

Sara pulled down her sleeve, which had pushed itself up.

<I always was a thin child,> she said bravely, <and I always had big green eyes.>

<I love your queer eyes,> said Ermengarde, looking into them with affectionate admiration. <They always look as if they saw such a long way. I love them—and I love them to be green—though they look black generally.>

<They are cat's eyes,> laughed Sara; <but I can't see in the dark with them—because I have tried, and I couldn't—I wish I could.>

It was just at this minute that something happened at the skylight which neither of them saw. If either of them had chanced to turn and look, she would have been startled by the sight of a dark face which peered cautiously into the room and disappeared as quickly and almost as silently as it had appeared. Not \textsc{quite} as silently, however. Sara, who had keen ears, suddenly turned a little and looked up at the roof.

<That didn't sound like Melchisedec,> she said. <It wasn't scratchy enough.>

<What?> said Ermengarde, a little startled.

<Didn't you think you heard something?> asked Sara.

<N-no,> Ermengarde faltered. <Did you?> {another ed. has <No-no,>}

<Perhaps I didn't,> said Sara; <but I thought I did. It sounded as if something was on the slates—something that dragged softly.>

<What could it be?> said Ermengarde. <Could it be—robbers?>

<No,> Sara began cheerfully. <There is nothing to steal\longdash>

She broke off in the middle of her words. They both heard the sound that checked her. It was not on the slates, but on the stairs below, and it was Miss Minchin's angry voice. Sara sprang off the bed, and put out the candle.

<She is scolding Becky,> she whispered, as she stood in the darkness. <She is making her cry.>

<Will she come in here?> Ermengarde whispered back, panic-stricken.

<No. She will think I am in bed. Don't stir.>

It was very seldom that Miss Minchin mounted the last flight of stairs. Sara could only remember that she had done it once before. But now she was angry enough to be coming at least part of the way up, and it sounded as if she was driving Becky before her.

<You impudent, dishonest child!> they heard her say. <Cook tells me she has missed things repeatedly.>

<'T warn't me, mum,> said Becky sobbing. <I was 'ungry enough, but 't warn't me—never!>

<You deserve to be sent to prison,> said Miss Minchin's voice. <Picking and stealing! Half a meat pie, indeed!>

<'T warn't me,> wept Becky. <I could 'ave eat a whole un—but I never laid a finger on it.>

Miss Minchin was out of breath between temper and mounting the stairs. The meat pie had been intended for her special late supper. It became apparent that she boxed Becky's ears.

<Don't tell falsehoods,> she said. <Go to your room this instant.>

Both Sara and Ermengarde heard the slap, and then heard Becky run in her slipshod shoes up the stairs and into her attic. They heard her door shut, and knew that she threw herself upon her bed.

<I could 'ave e't two of 'em,> they heard her cry into her pillow. <An' I never took a bite. 'Twas cook give it to her policeman.>

Sara stood in the middle of the room in the darkness. She was clenching her little teeth and opening and shutting fiercely her outstretched hands. She could scarcely stand still, but she dared not move until Miss Minchin had gone down the stairs and all was still.

<The wicked, cruel thing!> she burst forth. <The cook takes things herself and then says Becky steals them. She \textsc{doesn't}! She \textsc{doesn't}! She's so hungry sometimes that she eats crusts out of the ash barrel!> She pressed her hands hard against her face and burst into passionate little sobs, and Ermengarde, hearing this unusual thing, was overawed by it. Sara was crying! The unconquerable Sara! It seemed to denote something new—some mood she had never known. Suppose—suppose—a new dread possibility presented itself to her kind, slow, little mind all at once. She crept off the bed in the dark and found her way to the table where the candle stood. She struck a match and lit the candle. When she had lighted it, she bent forward and looked at Sara, with her new thought growing to definite fear in her eyes.

<Sara,> she said in a timid, almost awe-stricken voice, <are—are—you never told me—I don't want to be rude, but—are \textsc{you} ever hungry?>

It was too much just at that moment. The barrier broke down. Sara lifted her face from her hands.

<Yes,> she said in a new passionate way. <Yes, I am. I'm so hungry now that I could almost eat you. And it makes it worse to hear poor Becky. She's hungrier than I am.>

Ermengarde gasped.

<Oh, oh!> she cried woefully. <And I never knew!>

<I didn't want you to know,> Sara said. <It would have made me feel like a street beggar. I know I look like a street beggar.>

<No, you don't—you don't!> Ermengarde broke in. <Your clothes are a little queer—but you couldn't look like a street beggar. You haven't a street-beggar face.>

<A little boy once gave me a sixpence for charity,> said Sara, with a short little laugh in spite of herself. <Here it is.> And she pulled out the thin ribbon from her neck. <He wouldn't have given me his Christmas sixpence if I hadn't looked as if I needed it.>

Somehow the sight of the dear little sixpence was good for both of them. It made them laugh a little, though they both had tears in their eyes.

<Who was he?> asked Ermengarde, looking at it quite as if it had not been a mere ordinary silver sixpence.

<He was a darling little thing going to a party,> said Sara. <He was one of the Large Family, the little one with the round legs—the one I call Guy Clarence. I suppose his nursery was crammed with Christmas presents and hampers full of cakes and things, and he could see I had nothing.>

Ermengarde gave a little jump backward. The last sentences had recalled something to her troubled mind and given her a sudden inspiration.

<Oh, Sara!> she cried. <What a silly thing I am not to have thought of it!>

<Of what?>

<Something splendid!> said Ermengarde, in an excited hurry. <This very afternoon my nicest aunt sent me a box. It is full of good things. I never touched it, I had so much pudding at dinner, and I was so bothered about papa's books.> Her words began to tumble over each other. <It's got cake in it, and little meat pies, and jam tarts and buns, and oranges and red-currant wine, and figs and chocolate. I'll creep back to my room and get it this minute, and we'll eat it now.>

Sara almost reeled. When one is faint with hunger the mention of food has sometimes a curious effect. She clutched Ermengarde's arm.

<Do you think—you \textsc{could}?> she ejaculated.

<I know I could,> answered Ermengarde, and she ran to the door—opened it softly—put her head out into the darkness, and listened. Then she went back to Sara. <The lights are out. Everybody's in bed. I can creep—and creep—and no one will hear.>

It was so delightful that they caught each other's hands and a sudden light sprang into Sara's eyes.

<Ermie!> she said. <Let us \textsc{pretend}! Let us pretend it's a party! And oh, won't you invite the prisoner in the next cell?>

<Yes! Yes! Let us knock on the wall now. The jailer won't hear.>

Sara went to the wall. Through it she could hear poor Becky crying more softly. She knocked four times.

<That means, <Come to me through the secret passage under the wall,> she explained. <I have something to communicate.>>

Five quick knocks answered her.

<She is coming,> she said.

Almost immediately the door of the attic opened and Becky appeared. Her eyes were red and her cap was sliding off, and when she caught sight of Ermengarde she began to rub her face nervously with her apron.

<Don't mind me a bit, Becky!> cried Ermengarde.

<Miss Ermengarde has asked you to come in,> said Sara, <because she is going to bring a box of good things up here to us.>

Becky's cap almost fell off entirely, she broke in with such excitement.

<To eat, miss?> she said. <Things that's good to eat?>

<Yes,> answered Sara, <and we are going to pretend a party.>

<And you shall have as much as you \textsc{want} to eat,> put in Ermengarde. <I'll go this minute!>

She was in such haste that as she tiptoed out of the attic she dropped her red shawl and did not know it had fallen. No one saw it for a minute or so. Becky was too much overpowered by the good luck which had befallen her.

<Oh, miss! oh, miss!> she gasped; <I know it was you that asked her to let me come. It—it makes me cry to think of it.> And she went to Sara's side and stood and looked at her worshipingly.

But in Sara's hungry eyes the old light had begun to glow and transform her world for her. Here in the attic—with the cold night outside—with the afternoon in the sloppy streets barely passed—with the memory of the awful unfed look in the beggar child's eyes not yet faded—this simple, cheerful thing had happened like a thing of magic.

She caught her breath.

<Somehow, something always happens,> she cried, <just before things get to the very worst. It is as if the Magic did it. If I could only just remember that always. The worst thing never \textsc{quite} comes.>

She gave Becky a little cheerful shake.

<No, no! You mustn't cry!> she said. <We must make haste and set the table.>

<Set the table, miss?> said Becky, gazing round the room. <What'll we set it with?>

Sara looked round the attic, too.

<There doesn't seem to be much,> she answered, half laughing.

That moment she saw something and pounced upon it. It was Ermengarde's red shawl which lay upon the floor.

<Here's the shawl,> she cried. <I know she won't mind it. It will make such a nice red tablecloth.>

They pulled the old table forward, and threw the shawl over it. Red is a wonderfully kind and comfortable colour. It began to make the room look furnished directly.

<How nice a red rug would look on the floor!> exclaimed Sara. <We must pretend there is one!>

Her eye swept the bare boards with a swift glance of admiration. The rug was laid down already.

<How soft and thick it is!> she said, with the little laugh which Becky knew the meaning of; and she raised and set her foot down again delicately, as if she felt something under it.

<Yes, miss,> answered Becky, watching her with serious rapture. She was always quite serious.

<What next, now?> said Sara, and she stood still and put her hands over her eyes. <Something will come if I think and wait a little>—in a soft, expectant voice. <The Magic will tell me.>

One of her favourite fancies was that on <the outside,> as she called it, thoughts were waiting for people to call them. Becky had seen her stand and wait many a time before, and knew that in a few seconds she would uncover an enlightened, laughing face.

In a moment she did.

<There!> she cried. <It has come! I know now! I must look among the things in the old trunk I had when I was a princess.>

She flew to its corner and kneeled down. It had not been put in the attic for her benefit, but because there was no room for it elsewhere. Nothing had been left in it but rubbish. But she knew she should find something. The Magic always arranged that kind of thing in one way or another.

In a corner lay a package so insignificant-looking that it had been overlooked, and when she herself had found it she had kept it as a relic. It contained a dozen small white handkerchiefs. She seized them joyfully and ran to the table. She began to arrange them upon the red table-cover, patting and coaxing them into shape with the narrow lace edge curling outward, her Magic working its spells for her as she did it.

<These are the plates,> she said. <They are golden plates. These are the richly embroidered napkins. Nuns worked them in convents in Spain.>

<Did they, miss?> breathed Becky, her very soul uplifted by the information.

<You must pretend it,> said Sara. <If you pretend it enough, you will see them.>

<Yes, miss,> said Becky; and as Sara returned to the trunk she devoted herself to the effort of accomplishing an end so much to be desired.

Sara turned suddenly to find her standing by the table, looking very queer indeed. She had shut her eyes, and was twisting her face in strange convulsive contortions, her hands hanging stiffly clenched at her sides. She looked as if she was trying to lift some enormous weight.

<What is the matter, Becky?> Sara cried. <What are you doing?>

Becky opened her eyes with a start.

<I was a-'pretendin',' miss,> she answered a little sheepishly; <I was tryin' to see it like you do. I almost did,> with a hopeful grin. <But it takes a lot o' stren'th.>

<Perhaps it does if you are not used to it,> said Sara, with friendly sympathy; <but you don't know how easy it is when you've done it often. I wouldn't try so hard just at first. It will come to you after a while. I'll just tell you what things are. Look at these.>

She held an old summer hat in her hand which she had fished out of the bottom of the trunk. There was a wreath of flowers on it. She pulled the wreath off.

<These are garlands for the feast,> she said grandly. <They fill all the air with perfume. There's a mug on the wash-stand, Becky. Oh—and bring the soap dish for a centrepiece.>

Becky handed them to her reverently.

<What are they now, miss?> she inquired. <You'd think they was made of crockery—but I know they ain't.>

<This is a carven flagon,> said Sara, arranging tendrils of the wreath about the mug. <And this>—bending tenderly over the soap dish and heaping it with roses—<is purest alabaster encrusted with gems.>

She touched the things gently, a happy smile hovering about her lips which made her look as if she were a creature in a dream.

<My, ain't it lovely!> whispered Becky.

<If we just had something for bonbon dishes,> Sara murmured. <There!>—darting to the trunk again. <I remember I saw something this minute.>

It was only a bundle of wool wrapped in red and white tissue paper, but the tissue paper was soon twisted into the form of little dishes, and was combined with the remaining flowers to ornament the candlestick which was to light the feast. Only the Magic could have made it more than an old table covered with a red shawl and set with rubbish from a long-unopened trunk. But Sara drew back and gazed at it, seeing wonders; and Becky, after staring in delight, spoke with bated breath.

<This 'ere,> she suggested, with a glance round the attic—<is it the Bastille now—or has it turned into somethin' different?>

<Oh, yes, yes!> said Sara. <Quite different. It is a banquet hall!>

<My eye, miss!> ejaculated Becky. <A blanket 'all!> and she turned to view the splendours about her with awed bewilderment.

<A banquet hall,> said Sara. <A vast chamber where feasts are given. It has a vaulted roof, and a minstrels' gallery, and a huge chimney filled with blazing oaken logs, and it is brilliant with waxen tapers twinkling on every side.>

<My eye, Miss Sara!> gasped Becky again.

Then the door opened, and Ermengarde came in, rather staggering under the weight of her hamper. She started back with an exclamation of joy. To enter from the chill darkness outside, and find one's self confronted by a totally unanticipated festal board, draped with red, adorned with white napery, and wreathed with flowers, was to feel that the preparations were brilliant indeed.

<Oh, Sara!> she cried out. <You are the cleverest girl I ever saw!>

<Isn't it nice?> said Sara. <They are things out of my old trunk. I asked my Magic, and it told me to go and look.>

<But oh, miss,> cried Becky, <wait till she's told you what they are! They ain't just—oh, miss, please tell her,> appealing to Sara.

So Sara told her, and because her Magic helped her she made her \textsc{almost} see it all: the golden platters—the vaulted spaces—the blazing logs—the twinkling waxen tapers. As the things were taken out of the hamper—the frosted cakes—the fruits—the bonbons and the wine—the feast became a splendid thing.

<It's like a real party!> cried Ermengarde.

<It's like a queen's table,> sighed Becky.

Then Ermengarde had a sudden brilliant thought.

<I'll tell you what, Sara,> she said. <Pretend you are a princess now and this is a royal feast.>

<But it's your feast,> said Sara; <you must be the princess, and we will be your maids of honour.>

<Oh, I can't,> said Ermengarde. <I'm too fat, and I don't know how. \textsc{you} be her.>

<Well, if you want me to,> said Sara.

But suddenly she thought of something else and ran to the rusty grate.

<There is a lot of paper and rubbish stuffed in here!> she exclaimed. <If we light it, there will be a bright blaze for a few minutes, and we shall feel as if it was a real fire.> She struck a match and lighted it up with a great specious glow which illuminated the room.

<By the time it stops blazing,> Sara said, <we shall forget about its not being real.>

She stood in the dancing glow and smiled.

<Doesn't it \textsc{look} real?> she said. <Now we will begin the party.>

She led the way to the table. She waved her hand graciously to Ermengarde and Becky. She was in the midst of her dream.

<Advance, fair damsels,> she said in her happy dream-voice, <and be seated at the banquet table. My noble father, the king, who is absent on a long journey, has commanded me to feast you.> She turned her head slightly toward the corner of the room. <What, ho, there, minstrels! Strike up with your viols and bassoons. Princesses,> she explained rapidly to Ermengarde and Becky, <always had minstrels to play at their feasts. Pretend there is a minstrel gallery up there in the corner. Now we will begin.>

They had barely had time to take their pieces of cake into their hands—not one of them had time to do more, when—they all three sprang to their feet and turned pale faces toward the door—listening—listening.

Someone was coming up the stairs. There was no mistake about it. Each of them recognized the angry, mounting tread and knew that the end of all things had come.

<It's—the missus!> choked Becky, and dropped her piece of cake upon the floor.

<Yes,> said Sara, her eyes growing shocked and large in her small white face. <Miss Minchin has found us out.>

Miss Minchin struck the door open with a blow of her hand. She was pale herself, but it was with rage. She looked from the frightened faces to the banquet table, and from the banquet table to the last flicker of the burnt paper in the grate.

<I have been suspecting something of this sort,> she exclaimed; <but I did not dream of such audacity. Lavinia was telling the truth.>

So they knew that it was Lavinia who had somehow guessed their secret and had betrayed them. Miss Minchin strode over to Becky and boxed her ears for a second time.

<You impudent creature!> she said. <You leave the house in the morning!>

Sara stood quite still, her eyes growing larger, her face paler. Ermengarde burst into tears.

<Oh, don't send her away,> she sobbed. <My aunt sent me the hamper. We're—only—having a party.>

<So I see,> said Miss Minchin, witheringly. <With the Princess Sara at the head of the table.> She turned fiercely on Sara. <It is your doing, I know,> she cried. <Ermengarde would never have thought of such a thing. You decorated the table, I suppose—with this rubbish.> She stamped her foot at Becky. <Go to your attic!> she commanded, and Becky stole away, her face hidden in her apron, her shoulders shaking.

Then it was Sara's turn again.

<I will attend to you tomorrow. You shall have neither breakfast, dinner, nor supper!>

<I have not had either dinner or supper today, Miss Minchin,> said Sara, rather faintly.

<Then all the better. You will have something to remember. Don't stand there. Put those things into the hamper again.>

She began to sweep them off the table into the hamper herself, and caught sight of Ermengarde's new books.

<And you>—to Ermengarde—<have brought your beautiful new books into this dirty attic. Take them up and go back to bed. You will stay there all day tomorrow, and I shall write to your papa. What would \textsc{he} say if he knew where you are tonight?>

Something she saw in Sara's grave, fixed gaze at this moment made her turn on her fiercely.

<What are you thinking of?> she demanded. <Why do you look at me like that?>

<I was wondering,> answered Sara, as she had answered that notable day in the schoolroom.

<What were you wondering?>

It was very like the scene in the schoolroom. There was no pertness in Sara's manner. It was only sad and quiet.

<I was wondering,> she said in a low voice, <what \textsc{my} papa would say if he knew where I am tonight.>

Miss Minchin was infuriated just as she had been before and her anger expressed itself, as before, in an intemperate fashion. She flew at her and shook her.

<You insolent, unmanageable child!> she cried. <How dare you! How dare you!>

She picked up the books, swept the rest of the feast back into the hamper in a jumbled heap, thrust it into Ermengarde's arms, and pushed her before her toward the door.

<I will leave you to wonder,> she said. <Go to bed this instant.> And she shut the door behind herself and poor stumbling Ermengarde, and left Sara standing quite alone.

The dream was quite at an end. The last spark had died out of the paper in the grate and left only black tinder; the table was left bare, the golden plates and richly embroidered napkins, and the garlands were transformed again into old handkerchiefs, scraps of red and white paper, and discarded artificial flowers all scattered on the floor; the minstrels in the minstrel gallery had stolen away, and the viols and bassoons were still. Emily was sitting with her back against the wall, staring very hard. Sara saw her, and went and picked her up with trembling hands.

<There isn't any banquet left, Emily,> she said. <And there isn't any princess. There is nothing left but the prisoners in the Bastille.> And she sat down and hid her face.

What would have happened if she had not hidden it just then, and if she had chanced to look up at the skylight at the wrong moment, I do not know—perhaps the end of this chapter might have been quite different—because if she had glanced at the skylight she would certainly have been startled by what she would have seen. She would have seen exactly the same face pressed against the glass and peering in at her as it had peered in earlier in the evening when she had been talking to Ermengarde.

But she did not look up. She sat with her little black head in her arms for some time. She always sat like that when she was trying to bear something in silence. Then she got up and went slowly to the bed.

<I can't pretend anything else—while I am awake,> she said. <There wouldn't be any use in trying. If I go to sleep, perhaps a dream will come and pretend for me.>

She suddenly felt so tired—perhaps through want of food—that she sat down on the edge of the bed quite weakly.

<Suppose there was a bright fire in the grate, with lots of little dancing flames,> she murmured. <Suppose there was a comfortable chair before it—and suppose there was a small table near, with a little hot—hot supper on it. And suppose>—as she drew the thin coverings over her—<suppose this was a beautiful soft bed, with fleecy blankets and large downy pillows. Suppose—suppose\longdash> And her very weariness was good to her, for her eyes closed and she fell fast asleep.


She did not know how long she slept. But she had been tired enough to sleep deeply and profoundly—too deeply and soundly to be disturbed by anything, even by the squeaks and scamperings of Melchisedec's entire family, if all his sons and daughters had chosen to come out of their hole to fight and tumble and play.

When she awakened it was rather suddenly, and she did not know that any particular thing had called her out of her sleep. The truth was, however, that it was a sound which had called her back—a real sound—the click of the skylight as it fell in closing after a lithe white figure which slipped through it and crouched down close by upon the slates of the roof—just near enough to see what happened in the attic, but not near enough to be seen.

At first she did not open her eyes. She felt too sleepy and—curiously enough—too warm and comfortable. She was so warm and comfortable, indeed, that she did not believe she was really awake. She never was as warm and cosy as this except in some lovely vision.

<What a nice dream!> she murmured. <I feel quite warm. I—don't—want—to—wake—up.>

Of course it was a dream. She felt as if warm, delightful bedclothes were heaped upon her. She could actually \textsc{feel} blankets, and when she put out her hand it touched something exactly like a satin-covered eider-down quilt. She must not awaken from this delight—she must be quite still and make it last.

But she could not—even though she kept her eyes closed tightly, she could not. Something was forcing her to awaken—something in the room. It was a sense of light, and a sound—the sound of a crackling, roaring little fire.

<Oh, I am awakening,> she said mournfully. <I can't help it—I can't.>

Her eyes opened in spite of herself. And then she actually smiled—for what she saw she had never seen in the attic before, and knew she never should see.

<Oh, I \textsc{haven't} awakened,> she whispered, daring to rise on her elbow and look all about her. <I am dreaming yet.> She knew it \textsc{must} be a dream, for if she were awake such things could not—could not be.

Do you wonder that she felt sure she had not come back to earth? This is what she saw. In the grate there was a glowing, blazing fire; on the hob was a little brass kettle hissing and boiling; spread upon the floor was a thick, warm crimson rug; before the fire a folding-chair, unfolded, and with cushions on it; by the chair a small folding-table, unfolded, covered with a white cloth, and upon it spread small covered dishes, a cup, a saucer, a teapot; on the bed were new warm coverings and a satin-covered down quilt; at the foot a curious wadded silk robe, a pair of quilted slippers, and some books. The room of her dream seemed changed into fairyland—and it was flooded with warm light, for a bright lamp stood on the table covered with a rosy shade.

She sat up, resting on her elbow, and her breathing came short and fast.

<It does not—melt away,> she panted. <Oh, I never had such a dream before.> She scarcely dared to stir; but at last she pushed the bedclothes aside, and put her feet on the floor with a rapturous smile.

<I am dreaming—I am getting out of bed,> she heard her own voice say; and then, as she stood up in the midst of it all, turning slowly from side to side—<I am dreaming it stays—real! I'm dreaming it \textsc{feels} real. It's bewitched—or I'm bewitched. I only \textsc{think} I see it all.> Her words began to hurry themselves. <If I can only keep on thinking it,> she cried, <I don't care! I don't care!>

She stood panting a moment longer, and then cried out again.

<Oh, it isn't true!> she said. <It \textsc{can't} be true! But oh, how true it seems!>

The blazing fire drew her to it, and she knelt down and held out her hands close to it—so close that the heat made her start back.

<A fire I only dreamed wouldn't be \textsc{hot},> she cried.

She sprang up, touched the table, the dishes, the rug; she went to the bed and touched the blankets. She took up the soft wadded dressing-gown, and suddenly clutched it to her breast and held it to her cheek.

<It's warm. It's soft!> she almost sobbed. <It's real. It must be!>

She threw it over her shoulders, and put her feet into the slippers.

<They are real, too. It's all real!> she cried. <I am \textsc{not}—I am \textsc{not} dreaming!>

She almost staggered to the books and opened the one which lay upon the top. Something was written on the flyleaf—just a few words, and they were these:

<To the little girl in the attic. From a friend.>

When she saw that—wasn't it a strange thing for her to do—she put her face down upon the page and burst into tears.

<I don't know who it is,> she said; <but somebody cares for me a little. I have a friend.>

She took her candle and stole out of her own room and into Becky's, and stood by her bedside.

<Becky, Becky!> she whispered as loudly as she dared. <Wake up!>

When Becky wakened, and she sat upright staring aghast, her face still smudged with traces of tears, beside her stood a little figure in a luxurious wadded robe of crimson silk. The face she saw was a shining, wonderful thing. The Princess Sara—as she remembered her—stood at her very bedside, holding a candle in her hand.

<Come,> she said. <Oh, Becky, come!>

Becky was too frightened to speak. She simply got up and followed her, with her mouth and eyes open, and without a word.

And when they crossed the threshold, Sara shut the door gently and drew her into the warm, glowing midst of things which made her brain reel and her hungry senses faint. <It's true! It's true!> she cried. <I've touched them all. They are as real as we are. The Magic has come and done it, Becky, while we were asleep—the Magic that won't let those worst things \textsc{ever} quite happen.>